\documentclass[]{report}
\usepackage{amsmath}
\usepackage[spanish]{babel}

% Título y detalles del TP
\title{Introducción al Modelado Continuo \\[0.5cm]
	Trabajo Práctico 1 \\[0.5cm]
	Trayectorias de Kepler y Corrección Relativista}

% Sección de integrantes y grupo
\author{
	\textbf{Grupo: 11} \\ [1cm]
	\begin{tabular}{llll}
		\textbf{Apellido} & \textbf{Nombre} & \textbf{Libreta} \\ \hline
		Wilders Azara & Santiago & 350/19  \\
		Peralta & Martín & 428/22 \\
		Vernay & Jerónimo & ??? \\
	\end{tabular}
}

%Fecha
\date{\today}

\begin{document}
	\maketitle
	%\begin{abstract}
	%\end{abstract}
	
	% --- Desarrollo de los ejercicios ---
	
	\section*{Ejercicio 1}
	
	$$\ddot{u}(\theta) + u(\theta) - \frac{1}{\alpha} - \delta u^2(\theta) = 0$$
	
	Reducción de orden:
	
	\begin{equation*}
		\begin{aligned}
			y_1 &= u \\
			y_2 &= \dot{u} \\
			\dot{y_2} &= \ddot{u}
		\end{aligned}
		\quad \implies \quad 
		\dot{y_2} + y_1 - 1/a - \delta y_1^2 = 0
		\quad \implies \quad
		\begin{aligned}
			\dot{y_1} &= y_2 \\
			\dot{y_2} &= \delta y_1^2 + 1/a - y_1
		\end{aligned}
	\end{equation*}
	
	El estado del sistema queda definido por la inversa de la distancia radial ($y_1 = u(\theta)$) y su tasa de cambio respecto al ángulo ($y_2 = \dot{u}(\theta)$). Como el ángulo $\theta$, la variable independiente, no aparece de forma explicita en las ecuaciones, obtenemos un sistema autónomo no lineal de primer orden.
	\newpage
	
	\section*{Ejercicio 2}
	
	\newpage
	\section*{Ejercicio 3}

	
\end{document}